% !PRINTED 116
\weberpage{116}
%<<CONT
and if all elements are of degree $1$ or $2$, the groups
\[
1, \alpha, \beta, \alpha\beta
\]
with the condition $\alpha\beta = \beta\alpha$. There are therefore only these two groups of degree $4$. If the elements are denoted by $1, \alpha, \beta, \gamma$, we have the two tables:
\[
\begin{array}{c|cccc}
 & 1 & \alpha & \beta & \gamma \\
\hline
1 & 1 & \alpha & \beta & \gamma \\
\alpha & \alpha & \beta & \gamma & 1 \\
\beta & \beta & \gamma & 1 & \alpha \\
\gamma & \gamma & 1 & \alpha & \beta
\end{array}
\qquad
\begin{array}{c|cccc}
 & 1 & \alpha & \beta & \gamma \\
\hline
1 & 1 & \alpha & \beta & \gamma \\
\alpha & \alpha & 1 & \gamma & \beta \\
\beta & \beta & \gamma & 1 & \alpha \\
\gamma & \gamma & \beta & \alpha & 1
\end{array}
\]

The case $n = 6$ can be completely settled in the following way: We denote the six elements by $1, \alpha, \beta, \gamma, \delta, \varepsilon$, so that 1 is the unit of the group. If among them there occurs an element of degree 6, the group is cyclic.

If we therefore set this case aside, the degrees of the elements $\alpha, \beta, \gamma, \delta, \varepsilon$ can only be 2 or 3. If $\alpha$ and $\beta$ are of degree $2$, then $\alpha\beta$ cannot be of degree $2$; for otherwise we should have
\[
\alpha = \alpha^{-1}, \quad \beta = \beta^{-1}, \quad \alpha\beta = \beta^{-1}\alpha^{-1} = \beta\alpha,
\]
and $1, \alpha, \beta, \alpha\beta$ would therefore be a group of degree $4$; but such a group cannot be a divisor of a group of degree $6$.

There must therefore occur at least one element of degree $3$, and if we denote such an element by $\alpha$, and $\gamma$ is not contained in $1, \alpha, \alpha^2$, the group arranges itself thus:
\weq{1}{S = 1, \alpha, \alpha^2, \gamma, \gamma\alpha, \gamma\alpha^2.}
To determine the condition that this be a group, we form
\[
\gamma S = \gamma, \gamma\alpha, \gamma\alpha^2, \gamma^2, \gamma^2\alpha, \gamma^2\alpha^2,
\]
which must be identical with $S$; and therefore
\[
\gamma^2 = 1 \text{ or } = \alpha \text{ or } = \alpha^2
\]
must hold. But if $\gamma^2 = \alpha$ or $= \alpha^2$, then $\gamma^3 = \alpha\gamma$ or $= \alpha^2\gamma$, and $\gamma^6$, but no lower power of $\gamma = 1$, i.\,e. $S$ is
%CONT>>
